\paragraph{Higman groups have no nontrivial locally finite tree action.}
Let
\[
 H=\langle a_0,a_1,a_2,a_3\mid
       a_{i+1}a_i a_{i+1}^{-1}=a_i^2\ (i\bmod4)\rangle.
\]
Every homomorphism from $H$ to the automorphism group of a locally
finite simplicial tree is trivial.

\begin{proof}
Subdivide edges once to remove inversions. A hyperbolic tree isometry
has positive translation length, and its square has twice that length.
Since each $a_i$ is conjugate to its square, every $a_i$ is elliptic.

We first show that each adjacent pair $a=a_i$, $b=a_{i+1}$ has a
common fixed vertex. Put $A=\operatorname{Fix}(a)$, a nonempty convex
subtree. The relation gives
\[
 A\subseteq\operatorname{Fix}(a^2)=bA.
\]
Choose $p\in\operatorname{Fix}(b)$ and let $z$ be the unique nearest
point of $A$ to $p$. The nearest point of $bA$ to $p$ is $bz$, and
$d(p,bz)=d(p,z)$. Since $z\in A\subseteq bA$, uniqueness of the
nearest point in $bA$ gives $bz=z$. Thus $z$ is fixed by both $a$ and $b$.
In the subdivided simplicial tree the point $z$ may be chosen a vertex.

An elliptic automorphism of a locally finite tree has finite orbits
on vertices: a fixed vertex centers finite invariant balls. Suppose
some $a_i$ is nontrivial, and let $p$ be the smallest prime dividing
the length of any nontrivial vertex orbit of any of the four generators.
Choose $i$ and a vertex orbit of $a_i$ whose length is divisible by $p$.
The adjacent pair $a_i,a_{i+1}$ fixes a vertex $z_i$. A sufficiently
large finite ball about $z_i$ contains this orbit and is invariant under
both generators. Let $\bar a,\bar b$ be their permutations of this ball,
and let $m$ be the order of $\bar a$. Then $p\mid m$ and
\[
 \bar b\bar a\bar b^{-1}=\bar a^2.
\]
Consequently $m$ is odd. The order of the automorphism
$\bar a\mapsto\bar a^2$ of the cyclic group of order $m$ divides
the order of $\bar b$. In particular the multiplicative order
$\operatorname{ord}_p(2)$ divides the order of $\bar b$.
Since $p$ is odd, $\operatorname{ord}_p(2)>1$ and it divides $p-1$.
Choose a prime $q$ dividing $\operatorname{ord}_p(2)$; then $q<p$.
The order of a finite permutation is the least common multiple of
its cycle lengths, so $\bar b$ has a vertex orbit whose length is
divisible by $q$. This contradicts the definition of $p$.
Thus all four generators act trivially.
\end{proof}

Local finiteness is essential to this proof. It does not claim that
$H$ has Serre's property FA: its familiar amalgam decomposition does
give a nontrivial action on a tree with infinite valence.
